\subsection*{把一次购买的完整历史接起来}

沿同一个 S 继续一条具体历史。Alice 的三次购买请求恢复同一 P1/C1/T1，A 的业务受理提交一次。平台调用发出两次，沿同一范围和 \code{pay-C1} 恢复 X1，资金效果一次，余额从 10000 分变为 7000 分。可信回调、查询与旧工作者的观察相继到达，首次成功 R 建立一份 O1、原结果与 EV1，并结束 T1，后来的观察恢复这些事实。

PA1 把 EV1 交给代理，确认丢失。PA2 再发同一 EV1，代理保存 M1、M2 两份副本。M1 第一次投递时，消费者共同提交 I1/F1/回报，ACK 丢失。M1 随后重投，M2 再第一次投递。履约处理器共被调用三次，后两次通过 I1 读回首次保存的结果，两份确认分别结束各自投递的责任。

在这条截至建单的历史中，实际发货次数为零，履约回报待办为一份。前面独立审计的历史导出例子另有输入来源，不加入当前队列交付次数。将来发货调度者开始运行，再沿其 dispatch 身份另计外部动作。表~\ref{tab:ch10-counts} 把当前历史中的不同对象放在相同观察范围内比较。

\begin{table}[htbp]
\centering\small
\caption{同一条历史中，尝试、投递、持久结果与外部效果分别计数}\label{tab:ch10-counts}
\begin{tabularx}{\linewidth}{@{}p{.31\linewidth}p{.13\linewidth}X@{}}
\toprule
观察对象 & 次数 & 持久判断依据\\
\midrule
购买请求尝试／新增受理 & 3／1 & 同一 S 对应的 P1/C1/T1\\
平台调用／资金效果 & 2／1 & 原范围、\code{pay-C1} 与 X1\\
新订单／付款成功事件 & 1／1 & R、唯一 O1 与成功事实对应 EV1\\
发布调用／代理副本 & 2／2 & PA1/PA2 与同一 EV1\\
履约交付／处理器调用 & 3／3 & M1 重投、M2 首投与接收历史\\
新增接收终态／新增初次履约单 & 1／1 & I1 与 B1 关联真实 F1\\
履约回报待办／实际发货效果 & 1／0 & 同 Q 保存的回报，尚无发货提交\\
\bottomrule
\end{tabularx}
\end{table}

这里的新增业务效果次数与数据库 COMMIT 的总次数还不同。恢复一次原结果可以开启、结束一个只读事务，保存冲突、安排重试、领取任务也分别提交。首次成功迁移新增 O、F 或资金效果，其他事务保存各自的观察、调度与处置进度。判断同一授权执行了几次，要沿该授权的效果身份核对实际结果。

\CTenFigure{ch10-original09}{S、付款、事件、副本与独立消费者效果沿各自身份恢复}{fig:ch10-chain}
\CTenFigure{ch10-counts}{一次购买的通信分叉与效果汇合：两次发布形成三次交付，Q 建立一份 F}{fig:ch10-counts}

\FloatBarrier
\subsection*{安全关系和继续机会}

各段的安全关系在相应提交位置建立。A 固定 S 与 P/C 和原输入，平台 E 按自己的身份约束建立 X。R 把 O、原结果和唯一成功事件共同保存，发布恢复保持 EV 与原快照，Q 把本地接收依据、真实 F 与回报共同保存。购买重试先按 S 找到原 P/C，支付恢复从 C 取得原平台范围和键，发布器从 EV 取得原快照，消费者按 I 读完整绑定，再沿 B 关联实际 F。每段的后来者都从本段已保存的数据取得下一步输入。

代理的发布受理、消费确认和 Kafka offset 都是特定通信进度。它们能指出哪一段已经转交责任，业务结果由对应数据库或外部资源证明。相同 I 的重复被恢复，另一 O2 或独立审计范围则进入自己的接收判断。相同 B 的不同事件又共同受到初次履约约束。

要继续完成，尚未完成的输入必须持续留在至少一个恢复位置：上游 Outbox、代理保留副本、履约库待办或有归属的核查。数据库最终恢复可用，持锁事务结束，消费者获得足够运行机会。消息目标持续有效，网络与代理提供再次交付机会，调度让旧工作最终得到处理。隔离项由负责者取得证据并作出决定，回报与发货各自需要自己的下游机会。

ACK 丢失可以增加交付和调用，Q 的原子提交使新增 F 保持为一。下游能力不足时，发货可能停在未知核查。换副本、清历史或恢复旧备份以后，这些已保存关系是否存在，将决定同一副本再次到达时系统还能否认出它。
